\subsection{SUBALGEBRAS. IDEALS. QUOTIENT ALGEBRAS}

Let A be a commutative ring and E an A-algebra. If F is a sub-A-module of W which is stable under the multiplication on E, the restriction to \(F \times F\) of the multiplication of E defines (with the A-module structure on F) an A-algebra structure on F. F, with this structure, is called a subalgebra of the A-algebra E. Every intersection of subalgebras of E is a subalgebra of E. For every family \((x_{i})_{i \in I}\) of elements of E, the intersection of the subalgebras of E containing all the \(x_{i}\) is called the subalgebra of E generated by the family \((x_{i})_{i \in I}\) and \((x_{i})_{i \in I}\) is called a generating system (or generating family) of this subalgebra. If \(u: E \to E'\) is an A-algebra homomorphism, the image \(u(F)\) of every subalgebra F of E is a subalgebra of \(E'\) .

Let E be an associative algebra. For every subset M of E, the set M' of elements of E which are permutable with all the elements of M is a subalgebra of E called the centralizer subalgebra of M in E (I, § 1, no. 5). The centralizer M'' of M' in E is also called the bicentralizer of M; clearly M ⊆ M''. It follows that M' is contained in its bicentralizer M'', which is just the centralizer of M''; but the relation M ⊆ M'' implies M'' ⊆ M', so that M' = M'' (cf.~Set Theory, III, § 1, no. 7, Proposition 2). If F is a subalgebra of E, the centre of F is the intersection F ∩ F' of F and its centralizer F' in E. Note that if

F is commutative, then \(F \subset F'\) and hence \(F' \supset F''\) ; the bicentralizer \(F''\) of F is in this case the centre of F.

For certain non-associative algebras (for example Lie algebras) the notions of centralizer of a subalgebra and centre are defined differently (Lie Groups and Lie algebras, I, § 1, no. 6).

A subset \(\mathfrak{a}\) of an A-algebra \(\mathbf{E}\) is called a left ideal (resp. right ideal) of \(\mathbf{E}\) when \(\mathfrak{a}\) is a sub-A-module of \(\mathbf{E}\) and the relations \(x \in \mathfrak{a}, y \in \mathbf{E}\) imply \(yx \in \mathfrak{a}\) (resp. \(xy \in \mathfrak{a}\) ). It amounts to the same to say that \(\mathfrak{a}\) is a left ideal of \(\mathbf{E}\) or a right ideal of the opposite algebra \(\mathbf{E}^0\) . A two-sided ideal of \(\mathbf{E}\) is a subset \(\mathfrak{a}\) of \(\mathbf{E}\) which is both a left ideal and a right ideal. When \(\mathbf{E}\) is associative and admits a unit element \(e\) , then, for \(\alpha \in \mathbf{A}\) and \(x \in \mathbf{E}\) , \(\alpha x = (\alpha e)x = x(\alpha e)\) by virtue of (2) (no. 1) and hence the (right, left, two-sided) ideals of the ring \(\mathbf{E}\) (I, §8, no. 6) are identical with the (right, left, two-sided) ideals of the algebra \(\mathbf{E}\) . Every sum and every intersection of left (resp. right, two-sided) ideals of the algebra \(\mathbf{E}\) is a left (resp. right, two-sided) ideal. The intersection of the left (resp. right, resp. two-sided) ideals containing a subset \(\mathbf{X}\) of \(\mathbf{E}\) is called the left (resp. right, resp. two-sided) ideal of \(\mathbf{E}\) generated by \(\mathbf{X}\) .

Let \(\mathfrak{b}\) be a two-sided ideal of an A-algebra E. If \(x \equiv x'\pmod{\mathfrak{b}}\) and \(y \equiv y'\pmod{\mathfrak{b}}\) , then

\[
x (y - y ^ {\prime}) \in \mathfrak {b} \quad \text { and } \quad (x - x ^ {\prime}) y ^ {\prime} \in \mathfrak {b}
\]

and hence \(xy \equiv x'y' \pmod{b}\) . Hence an internal law can be defined on the quotient A-module E/b, which is the quotient of the multiplication law \((x,y) \mapsto xy\) of E by the equivalence relation \(x \equiv x' \pmod{b}\) (I, § 1, no. 6). It is immediately verified that this quotient law is an A-bilinear mapping of \((\mathrm{E} / \mathfrak{b}) \times (\mathrm{E} / \mathfrak{b})\) into E/b; it therefore defines with the A-module structure on E/b an A-algebra structure on E/b. E/b, with this algebra structure, is called the quotient algebra of the algebra E by the two-sided ideal b. The canonical mapping \(p: \mathrm{E} \to \mathrm{E} / \mathfrak{b}\) is an algebra homomorphism.

Let E, E' be two A-algebras and \(u:E \rightarrow E'\) an algebra homomorphism. The image \(u(\mathbf{E})\) is a subalgebra of \(E'\) and the kernel \(b = u(0)\) is a two-sided ideal of E; further, in the canonical decomposition of u:

\[
\mathrm{E} \xrightarrow {p} \mathrm{E} / \mathfrak {b} \xrightarrow {v} u (\mathrm{E}) \xrightarrow {j} \mathrm{E} ^ {\prime}
\]

v is an algebra isomorphism. More generally, all the results of Chapter I, § 8, no. 9 are still valid (and also their proofs) when the word ``ring'' is replaced by ``algebra''.

Let A be a commutative ring and E an A-algebra. On the set \(\tilde{\mathbf{E}} = \mathbf{A} \times \mathbf{E}\) , we define the following laws of composition:

\[
\begin{array}{r} (\lambda , x) + (\mu , y) = (\lambda + \mu , x + y) \\ (\lambda , x) (\mu , y) = (\lambda \mu , x y + \mu x + \lambda y) \\ \lambda (\mu , x) = (\lambda \mu , \lambda x). \end{array}
\]

It is immediately verified that \(\tilde{E}\) , with these laws of composition, is an algebra over A and \((1,0)\) is a unit element of this algebra. The set \(\{0\}\times E\) is a two-sided ideal of \(\tilde{E}\) and \(x\mapsto(0,x)\) is an isomorphism of the algebra E onto the subalgebra \(\{0\}\times E\) , by means of which E and \(\{0\}\times E\) are identified. \(\tilde{E}\) is called the algebra derived from E by adjoining a unit element; it is associative (resp. commutative) if and only if E is.

